\chapter{Fonctions de référence}\label{ch:g10:reffunc}

Une poignée de \omterm{def:g10:functions:function}{fonctions} simples --- affines, carré, inverse, racine
carrée, cube --- apparaissent partout, seules ou combinées. Connaître
leurs \omterm{def:g10:functions:graph}{graphiques} et leurs variations par cœur transforme de nombreux
problèmes en un croquis rapide. Ce chapitre les étudie chacune à son
tour et démontre leurs variations avec la technique de comparaison du
\cref{ch:g10:functions}.

\section{Fonctions affines}

\begin{definition}[Fonction affine]\label{def:g10:reffunc:affine}
Une \emph{fonction affine}\index{fonction affine} est une \omterm{def:g10:functions:function}{fonction} de la
forme
\[
f(x) = mx + p,
\]
où $m$ et $p$ sont des \omterm{def:g10:numbers:sets}{nombres réels} fixes. Son \omterm{def:g10:functions:graph}{graphique} est une
droite~: $m$ est le \emph{coefficient directeur}\index{coefficient
directeur} et $p$ l'\emph{ordonnée à l'origine}\index{ordonnée à
l'origine} (la droite coupe l'axe vertical en $(0, p)$). Lorsque
$p = 0$ la \omterm{def:g10:functions:function}{fonction} est \emph{linéaire}~: $f(x) = mx$ exprime la
proportionnalité.
\end{definition}

\begin{proposition}[Coefficient directeur et variations]\label{prop:g10:reffunc:affine}
Soit $f(x) = mx + p$.
\begin{enumerate}
  \item Pour deux entrées distinctes $u \neq v$~:
        $m = \dfrac{f(v) - f(u)}{v - u}$ (le \omterm{def:g10:reffunc:affine}{coefficient directeur} est
        la variation de la sortie par unité de variation de l'entrée).
  \item Si $m > 0$, $f$ est strictement \omterm{def:g10:functions:variations}{croissante} sur $\R$~; si
        $m < 0$, strictement \omterm{def:g10:functions:variations}{décroissante}~; si $m = 0$, constante.
\end{enumerate}
\end{proposition}

\begin{proof}
1.~Calculer $f(v) - f(u) = (mv + p) - (mu + p) = m(v - u)$, puis
diviser par $v - u \neq 0$.

2.~Si $u < v$, alors $f(v) - f(u) = m(v-u)$ a le signe de $m$, car
$v - u > 0$~: $m$ positif donne $f(u) < f(v)$ (\omterm{def:g10:functions:variations}{croissante}), $m$ négatif
donne $f(u) > f(v)$ (\omterm{def:g10:functions:variations}{décroissante}).
\end{proof}

\begin{omfigure}
\begin{tikzpicture}
\begin{axis}[omaxis,
  width=7.8cm, height=5.6cm,
  xmin=-3.4, xmax=3.4, ymin=-2.8, ymax=4.4,
  xtick={-2,-1,1,2}, ytick={-1,1,2,3}]
  \addplot[omDef, thick, domain=-3:2.6] {0.5*x + 1};
  \addplot[omThm, thick, domain=-2.2:3] {-x + 1};
  \node[omDef, font=\small] at (axis cs:2.5,2.7) {$y = \tfrac12 x + 1$};
  \node[omThm, font=\small] at (axis cs:-2.0,2.6) {$y = -x + 1$};
  \draw[dashed] (axis cs:1,1.5) -- (axis cs:2,1.5) -- (axis cs:2,2);
  \node[font=\footnotesize, below] at (axis cs:1.5,1.5) {$1$};
  \node[font=\footnotesize, right] at (axis cs:2,1.75) {$m$};
\end{axis}
\end{tikzpicture}

{\small Deux \omterm{def:g10:reffunc:affine}{fonctions affines}~: \omterm{def:g10:functions:variations}{croissante} pour $m = \frac12 > 0$,
\omterm{def:g10:functions:variations}{décroissante} pour $m = -1 < 0$. Avancer d'une unité vers la droite
change $y$ de $m$.}
\end{omfigure}

\begin{example}\label{ex:g10:reffunc:affine}
Trouver la \omterm{def:g10:reffunc:affine}{fonction affine} dont le \omterm{def:g10:functions:graph}{graphique} passe par $A(1, 3)$ et
$B(4, 9)$. Le \omterm{def:g10:reffunc:affine}{coefficient directeur} est
\[
m = \frac{9 - 3}{4 - 1} = \frac{6}{3} = 2 .
\]
Puis $f(x) = 2x + p$, et $f(1) = 3$ donne $2 + p = 3$, donc $p = 1$~:
$f(x) = 2x + 1$. Vérification avec $B$~: $f(4) = 9$.
\end{example}

\section{La fonction carré}

\begin{proposition}[Fonction carré]\label{prop:g10:reffunc:square}
La \omterm{def:g10:functions:function}{fonction} $f(x) = x^2$, définie sur $\R$~:
\begin{enumerate}
  \item est strictement \omterm{def:g10:functions:variations}{décroissante} sur $\intoc{-\infty}{0}$ et
        strictement \omterm{def:g10:functions:variations}{croissante} sur $\intco{0}{+\infty}$, avec \omterm{def:g10:functions:extrema}{minimum} $0$
        en $x = 0$~;
  \item satisfait $f(-x) = f(x)$ pour tout $x$~: son \omterm{def:g10:functions:graph}{graphique}, une
        \emph{parabole}\index{parabole}, est symétrique par rapport à
        l'axe vertical.
\end{enumerate}
\end{proposition}

\begin{proof}
1.~Soit $0 \leq u < v$. Alors
\[
f(v) - f(u) = v^2 - u^2 = (v - u)(v + u) > 0,
\]
car $v - u > 0$ et $v + u > 0$~: $f$ est strictement \omterm{def:g10:functions:variations}{croissante} sur
$\intco{0}{+\infty}$. Si $u < v \leq 0$, alors $v - u > 0$ mais
$v + u < 0$, donc $f(v) - f(u) < 0$~: strictement \omterm{def:g10:functions:variations}{décroissante}. Enfin
$x^2 \geq 0 = f(0)$ pour tout $x$.

2.~$(-x)^2 = x^2$~; donc les points $(x, x^2)$ et $(-x, x^2)$, \omterm{def:g10:functions:function}{images}
l'un de l'autre par symétrie par rapport à l'axe vertical, sont tous
deux sur le \omterm{def:g10:functions:graph}{graphique}.
\end{proof}

\begin{remark}[Carrés et ordre]
Parce que la \omterm{def:g10:functions:function}{fonction} carré décroît du côté négatif, élever au carré
\emph{inverse} l'ordre des nombres négatifs~: $-3 < -2$ mais $9 > 4$.
Ne jamais élever au carré les deux membres d'une inégalité sans vérifier
les signes.
\end{remark}

\section{La fonction inverse}

\begin{proposition}[Fonction inverse]\label{prop:g10:reffunc:inverse}
La \omterm{def:g10:functions:function}{fonction} $f(x) = \dfrac1x$, définie pour $x \neq 0$~:
\begin{enumerate}
  \item est strictement \omterm{def:g10:functions:variations}{décroissante} sur $\intoo{-\infty}{0}$ et
        strictement \omterm{def:g10:functions:variations}{décroissante} sur $\intoo{0}{+\infty}$~;
  \item satisfait $f(-x) = -f(x)$~: son \omterm{def:g10:functions:graph}{graphique}, une
        \emph{hyperbole}\index{hyperbole}, est symétrique par rapport à
        l'origine.
\end{enumerate}
\end{proposition}

\begin{proof}
1.~Soit $0 < u < v$. Alors
\[
f(v) - f(u) = \frac1v - \frac1u = \frac{u - v}{uv} < 0,
\]
car $u - v < 0$ et $uv > 0$~: strictement \omterm{def:g10:functions:variations}{décroissante} sur
$\intoo{0}{+\infty}$. Sur $\intoo{-\infty}{0}$ le même quotient a
$u - v < 0$ et $uv > 0$ encore (produit de deux négatifs), donc $f$
décroît aussi là.

2.~$\frac{1}{-x} = -\frac1x$.
\end{proof}

\begin{remark}
Attention~: $\frac1x$ n'est \emph{pas} \omterm{def:g10:functions:variations}{décroissante} sur tout son
\omterm{def:g10:functions:function}{ensemble de définition}. De $u = -1$ à $v = 1$ la valeur saute de $-1$ à
$1$. Les deux branches doivent être étudiées séparément.
\end{remark}

\begin{omfigure}
\begin{tikzpicture}
\begin{axis}[omaxis,
  width=5.6cm, height=5.2cm,
  xmin=-3.2, xmax=3.2, ymin=-2.6, ymax=6.4,
  xtick={-2,-1,1,2}, ytick={1,4},
  title={$y = x^2$}]
  \addplot[omDef, thick, domain=-2.45:2.45, samples=90] {x^2};
\end{axis}
\end{tikzpicture}\hfill
\begin{tikzpicture}
\begin{axis}[omaxis,
  width=5.6cm, height=5.2cm,
  xmin=-3.4, xmax=3.4, ymin=-3.4, ymax=3.4,
  xtick={-1,1}, ytick={-1,1},
  title={$y = \frac1x$}]
  \addplot[omDef, thick, domain=0.3:3.1, samples=80] {1/x};
  \addplot[omDef, thick, domain=-3.1:-0.3, samples=80] {1/x};
\end{axis}
\end{tikzpicture}

{\small La parabole $y = x^2$ (symétrique par rapport à l'axe vertical)
et l'hyperbole $y = \frac1x$ (deux branches, symétrique par rapport à
l'origine).}
\end{omfigure}

\section{Racine carrée et cube}

\begin{proposition}[Fonction racine carrée]\label{prop:g10:reffunc:sqrt}
La \omterm{def:g10:functions:function}{fonction} $f(x) = \sqrt{x}$, définie sur $\intco{0}{+\infty}$, est
strictement \omterm{def:g10:functions:variations}{croissante}.
\end{proposition}

\begin{proof}
Soit $0 \leq u < v$. Multiplier et diviser par le
\emph{conjugué}~:
\[
\sqrt v - \sqrt u
= \frac{(\sqrt v - \sqrt u)(\sqrt v + \sqrt u)}{\sqrt v + \sqrt u}
= \frac{v - u}{\sqrt v + \sqrt u} > 0,
\]
car $v - u > 0$ et $\sqrt v + \sqrt u > 0$ (noter $v > 0$, donc
$\sqrt v > 0$).
\end{proof}

\begin{proposition}[Fonction cube]\label{prop:g10:reffunc:cube}
La \omterm{def:g10:functions:function}{fonction} $f(x) = x^3$, définie sur $\R$, est strictement \omterm{def:g10:functions:variations}{croissante},
et son \omterm{def:g10:functions:graph}{graphique} est symétrique par rapport à l'origine.
\end{proposition}
\admitted

\begin{omfigure}
\begin{tikzpicture}
\begin{axis}[omaxis,
  width=5.6cm, height=5.0cm,
  xmin=-0.7, xmax=5.4, ymin=-0.7, ymax=2.9,
  xtick={1,4}, ytick={1,2},
  title={$y = \sqrt x$}]
  \addplot[omDef, thick, domain=0:5.1, samples=100] {sqrt(x)};
\end{axis}
\end{tikzpicture}\hfill
\begin{tikzpicture}
\begin{axis}[omaxis,
  width=5.6cm, height=5.0cm,
  xmin=-2.2, xmax=2.2, ymin=-3.4, ymax=3.4,
  xtick={-1,1}, ytick={-1,1},
  title={$y = x^3$}]
  \addplot[omDef, thick, domain=-1.48:1.48, samples=90] {x^3};
\end{axis}
\end{tikzpicture}

{\small La racine carrée (définie pour $x \geq 0$ seulement) et le cube,
toutes deux \omterm{def:g10:functions:variations}{croissantes}.}
\end{omfigure}

\begin{proposition}[Comparer $x$, $x^2$ et $\sqrt x$]\label{prop:g10:reffunc:compare}
Pour $0 < x < 1$~:\ \ $x^2 < x < \sqrt x$.\quad
Pour $x > 1$~:\ \ $\sqrt x < x < x^2$.\quad
En $x = 0$ et $x = 1$ les trois valeurs sont égales.
\end{proposition}

\begin{proof}
Supposons $0 < x < 1$. Multiplier $x < 1$ par $x > 0$ donne $x^2 < x$.
Ensuite, $x < \sqrt x$~: les deux membres étant positifs et le carré
croissant sur les positifs, cette inégalité équivaut à $x^2 < x$, que
l'on vient de prouver. Pour $x > 1$, multiplier $x > 1$ par $x$ donne
$x^2 > x$, et $\sqrt x < x$ suit par le même argument d'élévation au
carré.
\end{proof}

\begin{omfigure}
\begin{tikzpicture}
\begin{axis}[omaxis,
  width=7.4cm, height=5.8cm,
  xmin=-0.35, xmax=2.5, ymin=-0.35, ymax=2.9,
  xtick={1,2}, ytick={1,2}]
  \addplot[omThm, thick, domain=0:1.65, samples=80] {x^2};
  \addplot[omDef, thick, domain=0:2.35] {x};
  \addplot[omMeth, thick, domain=0:2.35, samples=80] {sqrt(x)};
  \fill (axis cs:1,1) circle (1.5pt);
  \node[omThm, font=\small, right] at (axis cs:1.45,2.4) {$y = x^2$};
  \node[omDef, font=\small, right] at (axis cs:2.0,1.85) {$y = x$};
  \node[omMeth, font=\small, right] at (axis cs:2.0,1.28) {$y = \sqrt x$};
\end{axis}
\end{tikzpicture}

{\small Les trois courbes se croisent en $(0,0)$ et $(1,1)$ et
échangent l'ordre là~: entre $0$ et $1$, $x^2$ est le plus petit et
$\sqrt x$ le plus grand~; au-delà de $1$ l'ordre est inversé.}
\end{omfigure}

\begin{method}[Comparer des images]\label{met:g10:reffunc:compare}
Pour comparer $f(a)$ et $f(b)$ pour une \omterm{def:g10:functions:function}{fonction} de référence $f$~:
\begin{enumerate}
  \item placer $a$ et $b$ sur un \omterm{def:g10:numbers:interval}{intervalle} où les variations de $f$
        sont connues~;
  \item si $f$ croît là, les \omterm{def:g10:functions:function}{images} sont dans le même ordre que les
        entrées~; si $f$ décroît, l'ordre est inversé~;
  \item si $a$ et $b$ ne sont pas dans le même \omterm{def:g10:numbers:interval}{intervalle} de monotonie,
        traiter d'abord les signes (par exemple pour le carré~: une
        entrée négative et une positive).
\end{enumerate}
\end{method}

\begin{example}
Comparer $\dfrac{1}{2.3}$ et $\dfrac{1}{2.4}$~: les deux entrées sont
dans $\intoo{0}{+\infty}$, où la \omterm{def:g10:functions:function}{fonction} inverse décroît, et
$2.3 < 2.4$, donc $\dfrac{1}{2.3} > \dfrac{1}{2.4}$. Comparer
$(-1.2)^2$ et $(-1.5)^2$~: sur $\intoc{-\infty}{0}$ le carré décroît, et
$-1.5 < -1.2$, donc $(-1.5)^2 > (-1.2)^2$ --- en effet $2.25 > 1.44$.
\end{example}

\section{Exercices}

\begin{exercise}[$\star$]\label{exo:g10:reffunc:1}
Pour chaque \omterm{def:g10:reffunc:affine}{fonction affine}, donner son \omterm{def:g10:reffunc:affine}{coefficient directeur}, son
\omterm{def:g10:reffunc:affine}{ordonnée à l'origine}, et dire si elle est \omterm{def:g10:functions:variations}{croissante} ou \omterm{def:g10:functions:variations}{décroissante}~:
\[
f(x) = 3x - 2, \qquad
g(x) = -\tfrac12 x + 4, \qquad
h(x) = 7, \qquad
k(x) = -x .
\]
\end{exercise}

\begin{exercise}[$\star$]\label{exo:g10:reffunc:2}
Trouver la \omterm{def:g10:reffunc:affine}{fonction affine} dont le \omterm{def:g10:functions:graph}{graphique} passe par $(2, 1)$ et
$(5, 10)$~; puis celle qui passe par $(-1, 4)$ et $(3, -4)$.
\end{exercise}

\begin{exercise}[$\star$]\label{exo:g10:reffunc:3}
Sans calculatrice, comparer~:
\[
(3.1)^2 \text{ et } (3.2)^2; \qquad
(-2.7)^2 \text{ et } (-2.8)^2; \qquad
\frac{1}{5.1} \text{ et } \frac{1}{5.2}; \qquad
\sqrt{17} \text{ et } \sqrt{15}.
\]
\end{exercise}

\begin{exercise}[$\star$]\label{exo:g10:reffunc:4}
En utilisant le \omterm{def:g10:functions:graph}{graphique} de la \omterm{def:g10:functions:function}{fonction} carré, résoudre $x^2 = 16$,
puis $x^2 \leq 16$, puis $x^2 > 9$.
\end{exercise}

\begin{exercise}[$\star$]\label{exo:g10:reffunc:5}
Soit $x \in \intoo{0}{1}$. Ordonner les nombres $x$, $x^2$, $x^3$ et
$\sqrt x$ du plus petit au plus grand, et vérifier avec $x = 0.25$.
\end{exercise}

\begin{exercise}[$\star\star$]\label{exo:g10:reffunc:6}
Résoudre graphiquement, puis algébriquement~: $\dfrac1x = 2$, et
$\dfrac1x < 2$ pour $x > 0$. Que change-t-on si l'on autorise aussi
$x < 0$~?
\end{exercise}

\begin{exercise}[$\star\star$]\label{exo:g10:reffunc:7}
Sachant que la \omterm{def:g10:functions:function}{fonction} carré décroît sur $\intoc{-\infty}{0}$ et croît
sur $\intco{0}{+\infty}$, encadrer $x^2$ lorsque~:
\[
\text{(a) } x \in \intcc{2}{5};
\qquad
\text{(b) } x \in \intcc{-3}{-1};
\qquad
\text{(c) } x \in \intcc{-2}{3}.
\]
(Dans le cas (c), attention~: le \omterm{def:g10:functions:extrema}{minimum} de $x^2$ n'est pas en une
extrémité.)
\end{exercise}

\begin{exercise}[$\star\star$]\label{exo:g10:reffunc:8}
Une compagnie de taxis facture un forfait fixe de $4$ plus $1.5$ par
kilomètre~; une autre ne facture pas de forfait et $2$ par kilomètre.
Modéliser chaque prix par une \omterm{def:g10:reffunc:affine}{fonction affine} de la distance $x$, tracer
les deux droites, et trouver à partir de quelle distance la première
compagnie est moins chère.
\end{exercise}

\begin{exercise}[$\star\star$]\label{exo:g10:reffunc:9}
Montrer que pour tous $a, b \geq 0$~:
$\sqrt{a + b} \leq \sqrt a + \sqrt b$. (Indication~: les deux membres
sont positifs ou nuls, donc comparer leurs carrés.) Quand y a-t-il
égalité~?
\end{exercise}

\begin{exercise}[$\star\star\star$]\label{exo:g10:reffunc:10}
Soit $f(x) = \dfrac{2x + 1}{x - 1}$, définie pour $x \neq 1$.
\begin{enumerate}
  \item Montrer que $f(x) = 2 + \dfrac{3}{x - 1}$ pour tout $x \neq 1$.
  \item En déduire les variations de $f$ sur $\intoo{1}{+\infty}$ à
        partir de celles de la \omterm{def:g10:functions:function}{fonction} inverse.
\end{enumerate}
\end{exercise}

\section{Problème~: les lois de la nature parlent en fonctions de
référence}

\begin{problem}[{Devoir du week-end --- les distances de freinage
croissent en $v^2$, les leviers obéissent à $1/d$, les planètes
suivent $a^{3/2}$~: lire la physique avec les \omterm{def:g10:functions:function}{fonctions} de
référence}]
\label{pb:g10:reffunc:1}
Ouvrez un livre de physique~: les mêmes quelques courbes reviennent à
chaque page --- la parabole, l'hyperbole, la courbe racine. La nature
écrit ses lois avec les \omterm{def:g10:functions:function}{fonctions} de référence de ce chapitre, et
connaître leurs allures (\cref{prop:g10:reffunc:square},
\cref{prop:g10:reffunc:inverse}, \cref{prop:g10:reffunc:sqrt},
\cref{prop:g10:reffunc:cube}) suffit à freiner une voiture, à
équilibrer un levier et à chronométrer les planètes. Le bouquet final
est la troisième loi de Kepler, lue directement sur les données du
système solaire.

\medskip
\noindent\textbf{Partie I --- La loi carrée du freinage.}
Une règle empirique pour la distance de freinage d'une voiture sur
route sèche~: $d = \left(\frac{v}{10}\right)^{2}$ mètres, à la vitesse
$v$ km/h.
\begin{enumerate}
  \item Calculer les distances de freinage à $30$, $50$, $90$ et
        $130$ km/h.
  \item Doubler la vitesse multiplie la distance de freinage par
        combien~? Expliquer à partir de l'effet d'échelle de la
        \omterm{def:g10:functions:function}{fonction} carré, et vérifier sur $50 \to 100$ km/h.
  \item Expertise judiciaire~: les traces de freinage mesurent
        $50$ m. De quelle vitesse la formule accuse-t-elle le
        conducteur (au km/h près)~? Quelle \omterm{def:g10:functions:function}{fonction} de référence a
        répondu --- et sur quel ensemble la lecture est-elle sans
        ambiguïté (\cref{prop:g10:reffunc:sqrt})~?
  \item Comparer le supplément de distance causé par \emph{un} km/h
        de plus, à basse et à haute vitesse~: calculer
        $d(31) - d(30)$ et $d(91) - d(90)$. Quelle propriété de la
        parabole (\cref{prop:g10:reffunc:square}) les deux réponses
        illustrent-elles~?
  \item Un arrêt réel ajoute le temps de réaction --- environ une
        seconde, pendant laquelle la voiture parcourt $0.28v$
        mètres~: $D(v) = 0.28v + \left(\frac{v}{10}\right)^2$.
        Calculer $D(50)$ et $D(130)$. Lequel des deux termes ---
        affine ou carré --- domine en ville, et lequel sur
        autoroute~?
\end{enumerate}

\medskip
\noindent\textbf{Partie II --- Les hyperboles du quotidien.}
\begin{enumerate}[resume]
  \item Un levier est en équilibre lorsque le produit force $\times$
        distance est le même des deux côtés. Un enfant de $60$ kg
        s'assied à $2$ m du pivot~; la force d'équilibre à la
        distance $d$ de l'autre côté vaut $F(d) = \frac{120}{d}$ (en
        équivalents-kg). Calculer $F(0.5)$, $F(1)$, $F(3)$.
  \item Que disent les variations de la \omterm{def:g10:functions:function}{fonction} inverse
        (\cref{prop:g10:reffunc:inverse}) à propos des leviers ---
        et pourquoi la fanfaronnade d'Archimède («~donnez-moi un
        point d'appui et je soulèverai la Terre~») se cache-t-elle
        dans la queue de l'hyperbole~? Qu'est-ce qui interdit
        $d = 0$~?
  \item Le son et la lumière s'affaiblissent comme le \emph{carré}
        de la distance~: à $1$ m d'une lampe l'intensité vaut $100$
        unités~; la donner à $2$, $3$ et $10$ m. Expliquer
        l'exposant $2$ à l'aide d'une sphère~: sur quelle aire la
        lumière s'est-elle répandue à la distance $d$ (problème de
        la pierre tombale d'Archimède, dans le volume précédent)~?
  \item Un car pour la sortie scolaire coûte $600$ euros, partagés
        également entre $n$ participants. Dresser un tableau du coût
        par tête pour $n = 10, 20, 30$, chercher combien de
        participants le ramènent à $25$ euros ou moins, et dire ce
        que l'allure de l'hyperbole promet --- et refuse --- quand
        $n$ grandit.
  \item Produire $x$ affiches coûte $2x + 3$ euros ($3$ euros de
        mise en route). Montrer que le coût \emph{moyen} par affiche
        vaut $a(x) = 2 + \frac3x$, décrire ses variations et
        interpréter économiquement son asymptote horizontale.
        (Comparer avec l'algèbre de l'\cref{exo:g10:reffunc:10}.)
\end{enumerate}

\medskip
\noindent\textbf{Partie III --- L'harmonie de Kepler.}
Pour chaque planète, on note $a$ sa distance moyenne au Soleil (en
unités astronomiques, Terre $= 1$) et $T$ sa période (en années).
Données~: Mars $a = 1.52$, $T = 1.88$~; Jupiter $a = 5.20$,
$T = 11.86$~; Saturne $a = 9.54$, $T = 29.4$.
\begin{enumerate}[resume]
  \item Calculer $a^3$ et $T^2$ pour la Terre, Mars et Jupiter.
        Qu'a remarqué Kepler en 1618~?
  \item Énoncer la loi comme une \omterm{def:g10:functions:function}{fonction}~: $T = \sqrt{a^3} =
        a\sqrt a$. La tester sur Saturne.
  \item Deux prédictions~: un astéroïde orbite à $a = 4$ UA ---
        trouver sa période (les nombres tombent juste)~; une comète
        revient tous les $27$ ans --- trouver sa distance moyenne
        (chercher un cube parfait).
  \item Quelles trois \omterm{def:g10:functions:function}{fonctions} de référence la loi tresse-t-elle
        ensemble~? Et sa règle d'échelle~: quand $a$ est multipliée
        par $4$, qu'advient-il de $T$~? (Vérifier avec l'astéroïde
        face à la Terre.)
  \item Kepler a lu cette loi dans les données de Tycho Brahe
        soixante-dix ans avant que la gravitation de Newton ne
        l'explique. En une phrase~: que dit cet épisode sur la
        puissance de reconnaître une \omterm{def:g10:functions:function}{fonction} de référence dans un
        tableau de nombres~?
\end{enumerate}

\medskip
\noindent\textbf{Partie IV --- Le lièvre et la tortue.}
\begin{enumerate}[resume]
  \item Ranger $\sqrt x$, $x$ et $x^2$ sur $\intoo{0}{1}$ puis sur
        $\intoo{1}{+\infty}$ (\cref{prop:g10:reffunc:compare}), et
        vérifier les deux classements en $x = 0.25$ et $x = 4$.
  \item Résoudre complètement $\sqrt x > x$ (pour $x \geq 0$, élever
        au carré en justifiant, puis conclure par un \omterm{def:g10:algebra:signtable}{tableau de
        signes}).
  \item Faire entrer le cube~: ranger $x$, $x^2$, $x^3$ en
        $x = 0.5$ puis en $x = 2$, et déterminer tous les points où
        les \omterm{def:g10:functions:function}{fonctions} carré et cube se croisent.
  \item Deux plans d'épargne rapportent, après $t$ années,
        $100\sqrt t$ euros (plan A) ou $10 t^2$ euros (plan B).
        Montrer que B dépasse A lorsque $t^3 = 100$, et donner
        l'instant de croisement au mois près. Morale sur les racines
        face aux puissances à long terme~?
  \item Pour finir~: pour chacune des lois de ce problème ---
        freinage ($v^2$), levier ($1/d$), lampe ($1/d^2$), Kepler
        ($a^{3/2}$), plan A ($\sqrt t$) --- dire en une proposition
        quelle particularité de sa courbe de référence porte le sens
        physique (montée qui se redresse, asymptote verticale,
        sphère qui s'étale, exposant d'échelle, croissance qui
        s'aplatit). Conclusion~: les \omterm{def:g10:functions:function}{fonctions} de référence sont
        l'alphabet~; la nature écrit avec.
\end{enumerate}
\end{problem}
